\usepackage{mathtext}
\usepackage{amsmath}
\usepackage{amsfonts}
\usepackage{amssymb}
\usepackage{latexsym}
\usepackage{mathtools}

\usepackage[russian]{babel}
\usepackage{cmap}
\usepackage{fontspec}
\setmainfont{Times New Roman}

\usepackage{graphicx}
\usepackage{csvsimple}
\usepackage{array}
\usepackage{longtable}
\usepackage{color}
\usepackage{listings}
\usepackage{pythonhighlight}
\usepackage{placeins}
\usepackage{pdfpages}
\usepackage{xurl}
\usepackage{csquotes}
\usepackage{enumitem}
\usepackage[top=20mm, right=10mm, left=29mm, bottom=20mm]{geometry}

% Поля и абзацы для TestVkr
\setlength{\parindent}{1.25cm}
\setlength{\parskip}{0pt}

% Дополнительно фиксируем красную строку после загрузки всех пакетов
\AtBeginDocument{%
    \setlength{\parindent}{1.25cm}%
    \setlength{\parskip}{0pt}%
}

\usepackage{setspace}
\onehalfspacing

\frenchspacing

\usepackage{indentfirst}
\usepackage{ragged2e}
\justifying

% Заголовки
\usepackage{titlesec}

\titleformat{name=\section}[block]
{\normalfont\Large\bfseries}
{\thesection}
{1em}
{}

\titleformat{name=\section,numberless}[block]
{\normalfont\Large\bfseries\centering}
{}
{0pt}
{}

\titleformat{name=\subsection}[block]
{\normalfont\normalsize\bfseries}
{\thesubsection}
{1em}
{}

\titleformat{name=\subsection,numberless}[block]
{\normalfont\normalsize\bfseries}
{}
{0pt}
{}

\titleformat{name=\subsubsection}[block]
{\normalfont\normalsize\bfseries}
{\thesubsubsection}
{1em}
{}

% Важно для TestVkr: номерные заголовки начинаются с красной строки
\titlespacing{\section}{\parindent}{*4}{*4}
\titlespacing{\subsection}{\parindent}{*4}{*4}
\titlespacing{\subsubsection}{\parindent}{*4}{*4}

\newcommand{\anonsection}[1]{%
    \section*{#1}%
    \addcontentsline{toc}{section}{#1}%
}

% Таблицы
\newenvironment{signstabular}[2][1]
{
    \renewcommand{\arraystretch}{#1}
    \begin{tabular}{#2}
        }
        {
    \end{tabular}
}

\newcommand{\specialcell}[2][c]{%
    \begin{tabular}[#1]{@{}c@{}}#2\end{tabular}%
}

% Графика
\graphicspath{{./}}
\DeclareGraphicsExtensions{.pdf,.png,.jpg,.eps,.svg}

% Подписи рисунков и таблиц
\usepackage[figurename=Рисунок]{caption}
\DeclareCaptionLabelSeparator{defffis}{~--~}

\captionsetup{
    justification=centering,
    singlelinecheck=false,
    labelsep=defffis
}

\captionsetup[table]{
    justification=raggedright,
    singlelinecheck=false,
    labelsep=defffis,
    skip=0pt,
    margin=1.25cm        % отступ слева и справа (или только слева)
}

% Список литературы
\usepackage[square,numbers,sort&compress]{natbib}
\renewcommand{\bibnumfmt}[1]{#1.\hfill}

\addto\captionsrussian{%
    \renewcommand{\refname}{Список использованных источников}%
}

\makeatletter
\renewenvironment{thebibliography}[1]
{
    \section*{\refname}%
    \@mkboth{\MakeUppercase\refname}{\MakeUppercase\refname}%
    \list{\@biblabel{\@arabic\c@enumiv}}%
    {
        \settowidth\labelwidth{\@biblabel{#1}}%
        \leftmargin=1.25cm
        \labelsep=0.5cm
        \@openbib@code
        \usecounter{enumiv}%
        \let\p@enumiv\@empty
        \renewcommand\theenumiv{\@arabic\c@enumiv}%
    }%
    \sloppy
    \clubpenalty4000
    \@clubpenalty \clubpenalty
    \widowpenalty4000%
    \sfcode`\.\@m
}
{
    \def\@noitemerr{%
        \@latex@warning{Empty `thebibliography' environment}%
    }%
    \endlist
}
\makeatother

% Дополнительные пакеты
\usepackage[figure,table]{totalcount}
\usepackage{ulem}
\usepackage{hhline}
\usepackage{rotating}
\usepackage{lastpage}

% Гиперссылки и перенос длинных URL
\usepackage[
    unicode=true,
    hidelinks,
    breaklinks=true
]{hyperref}

\Urlmuskip=0mu plus 1mu

% Списки
\renewcommand{\labelitemi}{---}

% Настройка нумерованных списков для TestVkr
\setlist[enumerate]{
    label=\arabic*.,
    leftmargin=0pt,
    itemindent=1.25cm,
    labelwidth=0.6cm,
    labelsep=0.2cm,
    align=left,
    itemsep=0.6em,
    topsep=0.6em,
    parsep=0pt
}

% Нумерация формул и рисунков по разделам
\numberwithin{equation}{section}
\numberwithin{figure}{section}

% Листинги программного кода
\renewcommand{\lstlistingname}{Листинг}

\lstdefinestyle{pythonstyle}{
    language=Python,
    basicstyle=\small\ttfamily,
    keywordstyle=\bfseries,
    commentstyle=\itshape,
    stringstyle=\ttfamily,
    numbers=left,
    numberstyle=\tiny,
    stepnumber=1,
    numbersep=8pt,
    frame=single,
    breaklines=true,
    breakatwhitespace=true,
    tabsize=4,
    showstringspaces=false,
    captionpos=b,
    columns=fullflexible,
    keepspaces=true
}

% ============================================
% НАСТРОЙКА СОДЕРЖАНИЯ
% ============================================

\usepackage{tocloft}

% Название содержания через babel
\addto\captionsrussian{%
    \renewcommand{\contentsname}{СОДЕРЖАНИЕ}%
}

% Заголовок содержания по центру
\renewcommand{\cfttoctitlefont}{\hfil\Large\bfseries}
\renewcommand{\cftaftertoctitle}{\hfil}
\setlength{\cftbeforetoctitleskip}{0pt}
\setlength{\cftaftertoctitleskip}{0.5cm}

% Точки от названия раздела до номера страницы
\renewcommand{\cftsecleader}{\cftdotfill{\cftdotsep}}
\renewcommand{\cftsubsecleader}{\cftdotfill{\cftdotsep}}
\renewcommand{\cftsubsubsecleader}{\cftdotfill{\cftdotsep}}

% Отступы в содержании для TestVkr
\cftsetindents{section}{0.50cm}{1.5em}
\cftsetindents{subsection}{1.55cm}{2em}
%\cftsetindents{subsubsection}{1.55cm}{2.5em}
\cftsetindents{subsubsection}{2.00cm}{2.5em}   % больше, чем subsection, но меньше 2.80

% ============================================
% ФИКС ОТСТУПОВ ДЛЯ TESTVKR
% ============================================

\AtBeginDocument{%
    \setlength{\parindent}{1.25cm}%
    \setlength{\parskip}{0pt}%
}

% ============================================
% ЖЕСТКАЯ ФИКСАЦИЯ АБЗАЦЕВ ДЛЯ TESTVKR
% ============================================

\setlength{\parindent}{1.25cm}
\setlength{\parskip}{0pt}

\AtBeginDocument{%
    \setlength{\parindent}{1.25cm}%
    \setlength{\parskip}{0pt}%
}

% Не даем окружениям сбрасывать красную строку
\makeatletter
\let\@afterindentfalse\@afterindenttrue
\@afterindenttrue
\makeatother